\documentclass{article}
\usepackage[localtheorems]{raeez-math-template}
\providecommand{\bC}{\mathbb{C}}
\providecommand{\bZ}{\mathbb{Z}}
\providecommand{\bQ}{\mathbb{Q}}
\providecommand{\Aut}{\mathrm{Aut}}
\providecommand{\Auts}{\mathrm{Aut}_s}
\providecommand{\Sp}{\mathrm{Sp}}
\providecommand{\Fix}{\mathrm{Fix}}
\providecommand{\Coh}{\mathrm{Coh}}
\providecommand{\CY}{\mathrm{CY}}
\providecommand{\DT}{\mathrm{DT}}
\providecommand{\red}{\mathrm{red}}
\providecommand{\ch}{\mathrm{ch}}
\providecommand{\BKM}{\mathrm{BKM}}
\providecommand{\cat}{\mathrm{cat}}
\providecommand{\Rep}{\mathrm{Rep}}
\providecommand{\fg}{\mathfrak{g}}
\providecommand{\kch}{\kappa_{\mathrm{ch}}}
\providecommand{\kBKM}{\kappa_{\mathrm{BKM}}}
\providecommand{\kcat}{\kappa_{\mathrm{cat}}}

\begin{document}

\section*{CHL action on the positive Hall half of $K3 \times E$}

\paragraph{The quotient.}
Fix $N \in \{1,2,3,4,6\}$ and a symplectic automorphism
$g_N \in \Auts(K3)$ of order $N$ (Nikulin, \emph{Finite automorphism
groups of K\"ahler K3 surfaces}, 1980). Fix an exact $N$-torsion
translation $t_N \in E[N]$. The CHL datum is the diagonal pair
$(g_N, t_N)$ acting diagonally on $K3 \times E$; the quotient
\[
  X_N \;=\; (K3 \times E)\big/\langle (g_N, t_N)\rangle
\]
is smooth, since $t_N$ acts freely on $E$ (Chaudhuri-Hockney-Lykken,
\emph{Maximally Supersymmetric String Theories in $D < 10$}, 1995,
arXiv:hep-th/9506144). Symplecticity of $g_N$ preserves the K3
holomorphic $2$-form; translation preserves the elliptic
$1$-form; the product $3$-form on $K3 \times E$ descends to $X_N$,
so $K_{X_N}\cong\mathcal{O}_{X_N}$.

The Hodge column is computed on invariants:
\[
  H^{0,\bullet}(X_N)
  \;\cong\;
  \bigl(H^{0,\bullet}(K3)\otimes H^{0,\bullet}(E)\bigr)^{
    \langle(g_N,t_N)\rangle}
  \;=\; (1,1,1,1).
\]
Thus
\[
  \kcat(X_N)=\chi(\mathcal{O}_{X_N})=1-1+1-1=0.
\]
The phrase CHL Calabi-Yau threefold records the descended holomorphic
volume form. In the strict algebraic convention $X_N$ is a smooth
$K$-trivial threefold with $h^{1,0}=h^{2,0}=1$.

\paragraph{The symplectic locus in the K3 moduli.}
Nikulin's theorem restricts $\Auts(K3)$ orders to divisors of elements
of $\{1,2,3,4,5,6,7,8\}$. The primitive CHL denominator slice uses
$N \in \{1,2,3,4,6\}$, the subset with
$\varphi(N) \mid 2$ (cyclotomic compatibility with paramodular
index $1$ in the CHL ladder). The orders $5,7,8$ belong to the
order-indexed Nikulin and Gritsenko-Clery automorphic rows, with a
separate row map.

The $g_N$-fixed sublattice
$\Lambda_{K3}^{g_N} \subset \Lambda_{K3} = \mathrm{II}_{3,19}$ is
computed by Nikulin (\emph{ibid.}, Table~1): of rank
$r_N \in \{22, 14, 10, 8, 6\}$ respectively (the invariant
dimensions of the induced representation on the $22$-dimensional
cohomology). The $g_N$-fixed K3 moduli sublocus is the
period domain
$\Omega(\Lambda_{K3}^{g_N}) \subset \Omega(\Lambda_{K3})$,
of dimension $r_N-2$ and codimension $22-r_N$ inside the full
$20$-dimensional K3 period domain. This is an orthogonal period
subdomain cut out by the anti-invariant lattice; Humbert divisors
enter after choosing the Borcherds product and its paramodular
realisation.

\paragraph{The $E[N]$-torsor on the elliptic factor.}
A point $t_N \in E[N]$ corresponds to the isogeny
$\pi_N \colon E \to E/\langle t_N\rangle$ of degree $N$.
The target $E/\langle t_N\rangle$ is $N$-isogenous to $E$. In the two
standard cyclic subgroups of $E_\tau=\bC/(\bZ+\tau\bZ)$ the quotient is
represented by $E_{N\tau}$ or $E_{\tau/N}$, and every cyclic subgroup of
order $N$ is obtained from these after a change of symplectic basis.
The locus
$\{t_N \in E[N] \mid t_N \text{ of exact order } N\}$ has
cardinality
\[
  N^2 \prod_{p\mid N}(1-p^{-2}).
\]
Over a ground field, descent asks that the pair $(g_N,t_N)$ and its
linearisation be Galois stable. All constructions here are over $\bC$.

\paragraph{The primitive denominator.}
Let $\Phi_N$ denote the primitive CHL Borcherds denominator attached to
the CHL-averaged Borcherds input. Its BKM character is
\[
  \operatorname{ch}(\fg_{\Phi_N})
  \;=\; \frac{1}{\Phi_N(Z)},
\]
and its weight is
\[
  k_N \;=\; \frac{c_N(0)}{2}
  \;\in\; \{5,\, 4,\, 3,\, 2,\, 1\}
  \quad\text{for } N \in \{1,2,3,4,6\}
\]
(Gritsenko, \emph{Arithmetical lifting and its applications}, 1999,
Theorem~1.2; Gritsenko-Nikulin~1995 Part~II Theorem~2.1). The primitive
CHL constants in this normalisation are
\[
  c_N(0)=(10,8,6,4,2).
\]
Singly twined K3 elliptic genera, fixed-locus traces, and the
Gritsenko-Clery eight-form catalogue carry their own constants. The
primitive CHL denominator has $k_6=1$.

Reduced curve counting uses an additional scalar convention. At
$N=1$, the Oberdieck-Pandharipande theorem for $K3\times E$ gives the
reduced Donaldson-Thomas scalar in the OP normalisation as the reciprocal
of a normalised square of $\Delta_5=\Phi_1$. The primitive BKM character
uses $1/\Phi_1$. For $N\geq2$, the Bryan-Oberdieck-Pandharipande-Yin CHL
curve-counting input and any comparison with the primitive denominator
requires specifying whether the scalar line is primitive or squared.
The primitive denominator convention is used throughout.

\paragraph{The positive Hall half.}
Let $Y^+_{\mathrm{Hall}}(K3\times E)$ denote the associative positive
Hall half of the oriented compact critical CoHA package, when that
package is fixed. For $d=3$ it is the Hall-level input whose
$\Phi_3$-image is $E_1$-chiral. The CHL generator acts on this positive
half by
\[
  \sigma_N \;=\; (g_N)_* \otimes (t_N)_* \colon\;
  Y^+_{\mathrm{Hall}}(K3\times E)
  \;\longrightarrow\;
  Y^+_{\mathrm{Hall}}(K3\times E),
\]
where $(g_N)_*$ is the induced automorphism on
$H^*_{T}(\mathrm{Hilb}(K3), \bQ)$ along the Maulik-Okounkov
$R$-matrix data. On ordinary cohomology $(t_N)_*$ is trivial; in the
quotient Hall problem it records the equivariant linearisation and the
torsion sector on the elliptic factor.

The quotient-stack positive half is the $\langle\sigma_N\rangle$-
equivariant Hall object. Its invariant decategorified algebra is
\[
  Y^+_{\mathrm{Hall}}(X_N)_{\mathrm{inv}}
  \;=\;
  \big(Y^+_{\mathrm{Hall}}(K3\times E)\big)^{\sigma_N},
\]
with charge grading landing in the $g_N$-invariant K3 Mukai lattice
together with the elliptic cohomology factor of
$E/\langle t_N\rangle$.

\paragraph{The Hall-Drinfeld boundary.}
Assume the compact Hall-Drinfeld recognition datum for the quotient:
an oriented compact critical CoHA, realised positive and negative
halves, a Cartan part, a Hopf pairing, radical quotient, parity
convention, PBW theorem, and Borcherds bracket comparison. Under these
hypotheses the completed Hall-Drinfeld double has Weyl-Kac-Borcherds
denominator $\Phi_N$ and
\[
  \kBKM(\Phi_N)=\frac{c_N(0)}2.
\]
This automorphic invariant is independent of the categorical Euler
characteristic
\[
  \kcat(X_N)=\chi(\mathcal{O}_{X_N})=0.
\]
The finite etale map $K3\times E\to X_N$ also gives
$\chi(\mathcal{O}_{K3\times E})=N\chi(\mathcal{O}_{X_N})$, hence the
same zero.

Since $d=3$, the CY-to-chiral construction supplies an $E_1$-chiral
output:
\[
  \Phi_3\bigl(D^b(\Coh(X_N))\bigr)\in E_1\text{-}\mathrm{ChirAlg}.
\]
The BKM superalgebra $\fg_{\Phi_N}$ is obtained after the K3-fibre
specialisation and the compact Hall-Drinfeld recognition datum above.
Its character is the primitive denominator character $1/\Phi_N(Z)$.

\paragraph{The boundary problem.}
At $N=1$, the automorphic denominator is the Gritsenko-Nikulin
$\Delta_5$ denominator and the reduced scalar is the
Oberdieck-Pandharipande square line. At $N\in\{2,3,4,6\}$, the
automorphic constants are fixed by
$c_N(0)=(8,6,4,2)$ and $\kBKM(\Phi_N)=(4,3,2,1)$. The remaining
mathematical construction is the compact equivariant Hall-Drinfeld
package whose positive half descends from
$Y^+_{\mathrm{Hall}}(K3\times E)$ and whose completed double has
denominator $\Phi_N$.

\end{document}
